\documentclass[a4paper,11pt]{article}
\usepackage[T2A]{fontenc}
\usepackage[utf8]{inputenc}
\usepackage[russian]{babel}
\usepackage{amsmath,amssymb,mathtools}
\usepackage{helvet}
\renewcommand{\familydefault}{\sfdefault}
\usepackage{icomma}
\usepackage[a4paper,left=21mm,right=21mm,top=21mm,bottom=22mm,headheight=14pt]{geometry}
\usepackage{microtype}
\usepackage{booktabs,tabularx,array}
\usepackage{enumitem}
\usepackage{xcolor}
\usepackage{fancyhdr}
\usepackage{tikz}
\usepackage{setspace}
\usepackage[hidelinks]{hyperref}
\definecolor{accent}{HTML}{007C80}
\definecolor{ink}{HTML}{24343A}
\definecolor{muted}{HTML}{63737A}
\setlength{\parindent}{0pt}
\setlength{\parskip}{0.30em}
\linespread{1.06}
\setlist[itemize]{leftmargin=*,nosep,topsep=3pt}
\pagestyle{fancy}
\fancyhf{}
\fancyhead[L]{\footnotesize\color{muted}Математика · Домашняя работа}
\fancyhead[R]{\footnotesize\color{muted}08.11.2026}
\fancyfoot[C]{\footnotesize\color{muted}\thepage}
\renewcommand{\headrulewidth}{0.3pt}
\renewcommand{\footrulewidth}{0pt}
\newcommand{\taskheading}[2]{%
  \phantomsection\label{task:#1}%
  {\color{accent}\large\bfseries Задание № #1}\par\vspace{2mm}%
  {\LARGE\bfseries\color{ink}#2}\par\vspace{3mm}\noindent\color{accent}\rule{\linewidth}{0.5pt}\color{black}\par}
\newcommand{\field}[1]{\par\textbf{#1}\par\nobreak}
\newcommand{\idea}{\textcolor{blue!60!black}{\textbf{Ключевая идея:}}}
\newcommand{\mistake}{\textcolor{red!75!black}{\textbf{Ошибка:}}}
\newcommand{\outcome}{\textcolor{teal!60!black}{\textbf{Итог:}}}
\begin{document}
\thispagestyle{empty}
\begin{center}
\vspace*{30mm}
{\small\bfseries\color{accent} МАТЕМАТИКА}\par
\vspace{18mm}
{\Huge\bfseries\color{ink}Ревью\\[4mm]домашней работы}\par
\vspace{10mm}
{\large Тригонометрические графики и прикладные задачи}\par
\vspace{22mm}
{\large Марк Гукин}\par
\vspace{13mm}
{\large 8 ноября 2026 года}\par
\vfill
{\normalsize Лёвин Артём Александрович}\par
{\small\color{muted}эксклюзивно для Марка Гукина}\par
\vspace{12mm}
\begin{tikzpicture}[x=1cm,y=1cm]
\draw[accent!65,line width=1pt] (-5.9,0) .. controls (-3.7,0.60) and (-1.8,-0.33) .. (0,0.10)
                               .. controls (2.3,0.59) and (3.9,-0.35) .. (5.9,0.03);
\end{tikzpicture}
\end{center}
\clearpage

{\LARGE\bfseries\color{ink}Сводка по заданиям, требующим разбора}\par
\medskip
В таблицу внесены задания, в которых обнаружены некорректные или недостаточно обоснованные шаги. Номера ведут к подробным разборам.\par
\vspace{6mm}
\renewcommand{\arraystretch}{1.38}
\small
\noindent\begin{tabularx}{\linewidth}{@{}>{\centering\arraybackslash}p{11mm}p{28mm}X>{\centering\arraybackslash}p{23mm}>{\centering\arraybackslash}p{25mm}@{}}
\toprule
\textbf{№} & \textbf{Тема} & \textbf{Суть замечания} & \textbf{Ваш ответ} & \textbf{Верный ответ} \\
\midrule
\hyperref[task:3]{3} & График синуса & При нахождении коэффициента пропущено слагаемое, хотя итог оказался верным. & $-4$ & $-4$ \\
\hyperref[task:7]{7} & Дифракционная решётка & В записи нарушено ограничение на угол; не обоснован выбор граничного значения. & $1800$ нм & $1800$ нм \\
\hyperref[task:10]{10} & Кинетическая энергия & При подстановке скорости пропущено возведение в квадрат. & $\sqrt{2}$ Дж & $4$ Дж \\
\bottomrule
\end{tabularx}
\normalsize
\vspace{10mm}
{\color{accent}\rule{\linewidth}{0.4pt}}\par
\textbf{На что обратить внимание.} Верный численный ответ ещё не гарантирует правильности всех переходов. При преобразовании формул полезно сохранять исходные слагаемые и степени, а в задачах с ограничением проверять, почему выбранное значение действительно даёт минимум или максимум.
\clearpage

\thispagestyle{empty}
\vspace*{0.32\textheight}
\begin{center}
{\fontsize{29}{36}\selectfont\bfseries\color{ink}Разбор ошибок}\par
\vspace{5mm}
{\color{accent}\rule{64mm}{0.65pt}}
\end{center}
\clearpage

\taskheading{3}{Синус и большой аргумент}
\field{Условие}
Функция $f(x)=a\sin x+b$ проходит через точки $A(0;-1)$, $B(\pi;-1)$ и $C\bigl(\frac{\pi}{2};2\bigr)$. Восстановите $a$ и $b$, затем вычислите $f\bigl(\frac{7\pi}{2}\bigr)$.

\field{Ваше решение и Ваш ответ}
Получено $b=-1$. После подстановки координат точки $C$ записано $2=a\cdot1$, затем сделан вывод $a=3$. Конечный ответ: $-4$.

\mistake\; При переходе к уравнению для $a$ потерян свободный член $b=-1$.

\field{Где именно возникает ошибка}
Последняя правильная запись при подстановке точки $C$:
\[
2=a\sin\frac{\pi}{2}+(-1).
\]
Следующая запись в тетради: $2=a\cdot1$. В ней уже нет слагаемого $-1$. Если принять её буквально, получилось бы $a=2$, что противоречит следующей строке $a=3$.

\field{Почему это неверно}
Значение синуса равно единице, поэтому заменяется только множитель $\sin\frac{\pi}{2}$. Свободный член $b$ при этом остаётся в уравнении. Вычитание единицы из правой части нельзя опускать.

\idea\; Разделяйте подстановку значений тригонометрической функции и перенос слагаемых: сперва вычислите синус, затем прибавьте единицу к обеим частям.

\field{Исправление и полное правильное решение}
По точке $A$:
\[
-1=a\sin0+b=0+b,\qquad b=-1.
\]
Точка $B$ подтверждает то же значение: $\sin\pi=0$; для нахождения $a$ нужна точка $C$. По ней:
\[
2=a\sin\frac{\pi}{2}-1=a-1,
\qquad a=2+1=3.
\]
Следовательно, $f(x)=3\sin x-1$. Приводим аргумент к стандартному углу:
\[
\frac{7\pi}{2}=2\pi+\frac{3\pi}{2},\qquad
\sin\frac{7\pi}{2}=\sin\frac{3\pi}{2}=-1.
\]
Тогда
\[
f\!\left(\frac{7\pi}{2}\right)=3\cdot(-1)-1=-4.
\]
\field{Проверка}
$f(0)=-1$, $f(\pi)=-1$, $f(\pi/2)=3-1=2$: все три точки принадлежат найденному графику.

\outcome\; Численный ответ $-4$ верный. Нужно исправить промежуточную запись: $2=a-1$.

\field{Задача на закрепление}
График функции $g(x)=c\sin x+d$ проходит через точки $(0;2)$, $(\pi;2)$ и $\bigl(\frac{\pi}{2};5\bigr)$. Найдите $g\bigl(\frac{7\pi}{2}\bigr)$.
\clearpage

\taskheading{7}{Минимальный период решётки}
\field{Условие}
Для света длиной волны $\lambda=300$ нм максимум третьего порядка наблюдается при $0^\circ<\varphi\le30^\circ$. По формуле $d\sin\varphi=k\lambda$ найдите минимальный период $d$; $k=3$.

\field{Ваше решение и Ваш ответ}
После записи ограничения $0<\varphi\le30^\circ$ появляется запись $\varphi>30^\circ$, противоречащая условию. Затем подставлено $\sin30^\circ=\frac12$, получено $d\cdot\frac12=900$ и $d=1800$ нм.

\mistake\; Граничное значение найдено верно, но направление ограничения на угол записано непоследовательно; не показано, почему именно при $30^\circ$ период наименьший.

\field{Где именно возникает ошибка}
Из условия следует $\varphi\le30^\circ$. Следующая строка $\varphi>30^\circ$ противоречит исходному ограничению. Подставить $30^\circ$ можно лишь после обоснования, что эта точка действительно даёт минимум $d$.

\field{Почему это неверно}
Из формулы
\[
d=\frac{k\lambda}{\sin\varphi}
\]
видно, что при положительном числителе большему значению синуса соответствует меньший период $d$. На разрешённом промежутке от $0^\circ$ до $30^\circ$ синус возрастает. Его наибольшее допустимое значение достигается в правой границе $30^\circ$.

\idea\; Для минимизации положительной дроби с постоянным числителем найдите максимально допустимое значение знаменателя.

\field{Исправление и полное правильное решение}
По формуле решётки:
\[
d\sin\varphi=3\cdot300=900\ \text{нм},
\qquad d=\frac{900}{\sin\varphi}\ \text{нм}.
\]
Поскольку $0^\circ<\varphi\le30^\circ$, то
\[
0<\sin\varphi\le\sin30^\circ=\frac12.
\]
Следовательно,
\[
d=\frac{900}{\sin\varphi}\ge\frac{900}{1/2}=1800\ \text{нм}.
\]
Равенство достигается при допустимом угле $\varphi=30^\circ$.

\field{Проверка}
$1800\cdot\sin30^\circ=1800\cdot\frac12=900=3\cdot300$. Значит, максимум третьего порядка возникает при требуемом угле.

\outcome\; Ответ $1800$ нм верный. Для строгого решения необходимо показать, что меньший период потребовал бы угол больше $30^\circ$.

\field{Задача на закрепление}
Для света с длиной волны $\lambda=250$ нм найдите минимальный период решётки, если максимум второго порядка должен наблюдаться при угле $0^\circ<\varphi\le30^\circ$. Используйте $d\sin\varphi=k\lambda$.
\clearpage

\taskheading{10}{Энергия гармонических колебаний}
\field{Условие}
Дано $m=1$ кг, $v_0=4$ м/с, $T=8$ с, $t=1$ с. Скорость и кинетическая энергия связаны формулами
\[
v(t)=v_0\sin\!\left(\frac{2\pi t}{T}\right),\qquad
E(t)=\frac{m\,[v(t)]^2}{2}.
\]
Найдите $E(1)$ в джоулях.

\field{Ваше решение и Ваш ответ}
Скорость вычислена верно: $v(1)=\dfrac{4\sqrt2}{2}=2\sqrt2$ м/с. В следующей части решения скорость подставлена в числитель формулы энергии без её возведения в квадрат, и получен результат $\sqrt2$ Дж.

\mistake\; Потеря квадрата скорости при переходе к формуле энергии.

\field{Где именно возникает ошибка}
Последний правильный шаг: $v(1)=2\sqrt2$ м/с. В следующем переходе вместо квадрата найденной скорости используется сама скорость. Правильный числитель содержит $\bigl(2\sqrt2\bigr)^2$.

\field{Почему это неверно}
В формуле $E=mv^2/2$ вторая степень относится ко всему значению скорости. Сначала нужно возвести в квадрат произведение $2\sqrt2$:
\[
(2\sqrt2)^2=2^2\cdot(\sqrt2)^2=4\cdot2=8.
\]
Без квадрата получается выражение, не соответствующее формуле кинетической энергии.

\idea\; Выписывайте численное значение скорости в скобках, если далее его нужно возводить в квадрат: $[v(1)]^2=(2\sqrt2)^2$.

\field{Исправление и полное правильное решение}
\[
\begin{aligned}
\frac{2\pi t}{T}&=\frac{2\pi\cdot1}{8}=\frac\pi4,\\
v(1)&=4\sin\frac\pi4=4\cdot\frac{\sqrt2}{2}=2\sqrt2\ \text{м/с},\\
E(1)&=\frac{1\cdot(2\sqrt2)^2}{2}=\frac{8}{2}=4\ \text{Дж}.
\end{aligned}
\]

\field{Проверка}
Максимальная скорость равна $4$ м/с, поэтому максимальная кинетическая энергия составляет $\frac{1\cdot4^2}{2}=8$ Дж. При скорости $2\sqrt2$ м/с энергия $4$ Дж; значение попадает в допустимый диапазон от нуля до $8$ Дж.

\outcome\; Правильная кинетическая энергия равна $4$ Дж.

\field{Задача на закрепление}
Груз массой $m=2$ кг колеблется со скоростью $v(t)=6\sin\!\left(\frac{2\pi t}{12}\right)$ м/с. Найдите кинетическую энергию при $t=2$ с, пользуясь формулой $E=mv^2/2$.
\vfill
\begin{center}
{\color{muted}\small Лёвин Артём Александрович эксклюзивно для Марка Гукина}
\end{center}
\end{document}
